\section{Typed Model Foundation and Claim Boundary}
\label{sec:oc133-integrated-scientific-closure}
The didactic opening has now given the reader an intuitive picture of a continuum and of the K-level hierarchy. This chapter turns that picture into the first scientific claim of the release. OC Core 1.3.3 treats a continuum as a typed model object: it has a carrier, admissible state region, boundary, axes of variation, thresholds, potentials, flows, cycles, continuumness, operators, and an embedding context. The claim is not that every domain has already been solved. The claim is that these components provide a common grammar for stating how a system persists, changes, fails, recovers, or becomes part of another system.

The public manuscript identifier is OC Core 1.3.3; its Concept DOI is printed on the title page. Platform events, repository state, and archival records are support surfaces rather than elements of the scientific argument. The argument itself is the chain from definition to intuition, from intuition to formal anchor, from formal anchor to proof or finite witness, and from evidence to an explicit limitation or falsifier.

\subsection{Claim, Intuition, and Formal Anchor}
The central intuition is that a living or coherent system is not exhausted by its parts. A system also has a region of admissible states, a boundary beyond which it no longer counts as the same live continuum, flows that move it, cycles that sustain it, and thresholds that can change its regime. In OC notation this is summarized by the tuple \textbackslash{}(K=(\textbackslash{}Omega,\textbackslash{}partial\textbackslash{}Omega,A,\textbackslash{}Theta,P,J,C,k,M)\textbackslash{}). The tuple is a compact scientific object, not a decorative diagram: each component gives a reviewer a place to ask whether the model has preserved a real distinction or merely renamed one.

A worked example can be read before the theorem material. Consider an AI-enabled enterprise service. Its lower continua include compute, data, model weights, prompts, monitoring, security controls, teams, contracts, and user trust. The service remains live only while its admissible technical, organizational, legal, and economic state remains inside a declared boundary. A model update can be technically successful and still push the enterprise continuum outside its admissible state if it breaks trust, governance, compliance, or operational resilience. OC's formal vocabulary exists to make that cross-level failure stateable.

\subsection{Evidence and Proof Anchors}
The release binds the model grammar to several evidence classes. The theorem route names formal statements and their assumptions. Proof sheets and the Lean subset record which parts are mechanized or proof-oriented. Finite semantic checks show bounded model behavior under explicit cases. Target-blind replay rows and numeric tables show how selected claims are reconstructed under fixed inputs, formulas, comparators, residuals, and negative controls. Prior-art and novelty material then ask what existing traditions already explain and what residual delta remains.

These anchors should be read in sequence. A theorem identifier is a citation handle, not a substitute for the argument. A finite witness is a bounded semantic case, not a proof that the declared release domains are complete. A replay row is evidence for the scoped reconstruction it names, not final bounded replay QA. A comparator row is a boundary against overclaiming, not a declaration that prior work is absent.

\subsection{Limitations and Falsifiers}
Every promoted claim remains reopenable. A claim can fail if its assumptions are contradicted, if the finite case no longer reproduces, if a negative control succeeds, if a stronger comparator explains the same case with less burden, if a proof dependency is lost, or if a domain example requires variables absent from the declared tuple. This is why the manuscript keeps limitations inside the scientific path. A limitation is not an apology; it is part of the meaning of the claim.

The rest of the monograph follows that pattern: claim, intuition, worked example, formal anchor, evidence or proof anchor, and limitation. The reader sees finished scientific prose. Machine-readable rows remain available for audit, but they no longer become chapter titles or narrative paragraphs.


\section{From Model Grammar to Scientific Reading}
\label{sec:oc133-reader-contract-scientific-route}
Different readers enter the model with different expectations. A formal-methods reader asks whether the definitions have assumptions and proof obligations. A systems-theory reader asks whether the model adds anything beyond earlier accounts of organization, feedback, and emergence. A domain reader asks whether the examples preserve the variables that matter in practice. The manuscript therefore develops the argument in layers rather than assuming one scoped starting point.

The first layer is conceptual: a continuum has admissible states, boundaries, thresholds, flows, cycles, and embedding context. The second layer is formal: those terms become typed objects, operators, witnesses, proof statements, and finite semantic cases. The third layer is evidential: selected claims are connected to replay rows, numeric anchors, comparator material, and negative controls. The fourth layer is critical: limitations and falsifiers state how the claim can be narrowed or reopened.

This layered reading is meant to protect both skepticism and usefulness. The theory should not be accepted because its vocabulary is broad; it should be tested where the vocabulary claims to preserve a distinction that ordinary domain language leaves implicit. The reader is invited to press exactly there: at boundaries, assumptions, examples, formulas, evidence classes, comparators, and reopening conditions.

The manuscript's practical promise is that a reader can move from an intuitive system example to a formal component and then to an evidence surface without changing the subject. If that movement breaks, the corresponding claim should be treated as unfinished.


\section{Literature and Prior-Art Position}
\label{sec:oc133-literature-prior-art-position}
The monograph carries 129 bibliography entries and source-backed comparator anchors with overlap and residual-delta notes. The literature layer locates OC against emergence, phase transitions, autopoiesis, dynamical systems, RAF chemistry, information and complexity measures, identity over time, hybrid systems, formal methods, systems engineering, and reproducibility practice. The release uses literature to bound novelty, not to claim absence of all prior work.

A scientific reviewer should therefore read the comparator material as a map of overlap and residual delta. Where OC integrates or rephrases an existing tradition, the monograph says so. Where a broader priority or unrestricted comparative claim is not evidenced, it stays outside the promoted release surface.


\section{How the Figures Carry the Argument}
\label{sec:oc133-figure-route-design-logic}
The full corpus includes 32 public figure entries across the theorem roadmap, worked examples, and inline evidence routes. The figure layer is part of the argument. A useful OC figure should show which variables matter, which boundary is being crossed or preserved, which formal expression it illustrates, and which evidence or falsifier would matter if the picture were wrong.

The early continuum and K-level figures therefore do more than decorate the opening pages. They show the state region, boundary, flow, threshold, cycle, continuumness condition, K0--K12 hierarchy, upward composition, and downward constraint before the notation becomes dense. Later figures should be read the same way: as compact demonstrations of a claim that the surrounding prose and evidence must still justify.


\section{Integrated 1.3.3 Scientific Argument Before Evidence Rows}
\label{sec:oc133-integrated-argument-before-evidence-rows}
Before the monograph names individual evidence rows, it states the integrated argument in ordinary scientific prose. OC Core 1.3.3 promotes a bounded model-core claim: a continuum is treated as a typed object whose carrier, realization, lawful possibility, liveness, residue, boundary behavior, operator behavior, cycle mode, dimension semantics, and K-level position can be inspected under declared assumptions. The release does not ask the reader to trust an inventory. It asks the reader to test whether those typed components solve specific ambiguity classes that appear when persistence, death, rebirth, boundary, update, and cross-level reduction are discussed without a common grammar.

The proof layer then gives the argument its ceiling. The promoted theorem surface is not a claim that OC has completed every possible domain science. It is a claim that named theorem statements have named assumptions, proof sheets, Lean-subset references where available, finite semantic witnesses, negative controls, and counterexample boundaries. A theorem identifier is therefore a citation handle, not the argument itself: the argument is the chain from definition to lemma to proof sheet to finite witness to explicit reopening condition.

The empirical layer has the same discipline. Physics, chemistry, biology, systems, and mathematics rows are target-blind or replay-QA examples with formulas, snapshots, comparators, residuals, negative controls, falsifiers, and hashes. They show how the OC release binds a claim to a replayable evidence route. They do not promote scoped numeric closure, final full-scope synthesis status, or unrestricted comparative claim over modern science. Those broader ambitions remain a background research program until they can be supported at the same or higher evidential standard.

The comparator layer prevents novelty theater. General systems theory, autopoiesis, dynamical systems, category and topos formalisms, RAF theory, complexity and information measures, identity theory, systems engineering, and reproducible-research practice are treated as live comparators. For each tradition, the release first records overlap, then states the residual delta, and finally names the claim that OC must not make. This order is essential: accepted overlap comes before residual contribution.

The row-based sections that follow are therefore not a substitute for the monograph. They are an audit trail. A reader who wants the conceptual argument should read the prose spine first; a reviewer who wants to attack exact wording can then use the rows to locate the proof sheet, finite case, replay row, comparator source, or reopening rule that carries the challenged sentence.
